% Authoritative results: upstream 4d6a4a2, sections/04_phase1_retrieval.tex.
\subsubsection{Kết quả và phân tích}
\begin{table}[H]\setstretch{1.05}
  \centering
  \caption{Kết quả sàng lọc retriever trên 281 câu resolved.}
  \label{tab:phase1-retriever-screening}
  \small
  \begin{tabular}{lrrrr}
    \toprule
    \textbf{Retriever} & \textbf{MRR@5} & \textbf{nDCG@5} &
    \textbf{Hit@5} & \textbf{P50 (ms)} \\
    \midrule
    BGE-M3 learned sparse          & \textbf{0,8059} & \textbf{0,8317} & \textbf{0,9181} & 76,7 \\
    BM25 Okapi + stemming          & 0,7815 & 0,8123 & 0,9146 & 56,1 \\
    BM25+                          & 0,6881 & 0,7206 & 0,8221 & 101,1 \\
    BM25 Okapi                     & 0,6732 & 0,7087 & 0,8185 & 111,8 \\
    Dense: e5-base-v2              & \textbf{0,6316} & \textbf{0,6684} & \textbf{0,7865} & 15,5 \\
    Dense: bge-small-en-v1.5       & 0,6285 & 0,6589 & 0,7651 & 16,0 \\
    Dense: bge-large-en-v1.5       & 0,6252 & 0,6536 & 0,7544 & 27,4 \\
    Dense: all-MiniLM-L6-v2        & 0,4851 & 0,5165 & 0,6299 & 10,6 \\
    \bottomrule
  \end{tabular}
\end{table}

Trên 281 câu resolved, BGE-M3 đạt điểm cao nhất trong nhóm sparse, còn E5-base-v2 được chọn làm đại diện dense. Các phương pháp sparse vượt trội rõ rệt so với các mô hình dense. Phân tích EDA chỉ ra rằng câu hỏi tin tức neo chặt vào các thực thể định danh rời rạc (tên riêng, địa danh, mốc thời gian với $\text{IDF} \ge 6{,}0$). Khi kho ngữ liệu lớn (11.064 bài), mô hình dense bi-encoder nén thông tin vào vector 768 chiều dễ gặp hiện tượng \textbf{pha loãng thực thể (entity dilution)}, làm mờ ranh giới giữa các sự kiện cùng chủ đề và kéo theo nhiều bài báo gây nhiễu; trong khi sparse dựa vào khớp từ vựng chính xác để cô lập không gian tìm kiếm. Tuy nhiên, CI95 của hiệu nDCG@5 giữa BGE-M3 và BM25 có stemming là $[-0{,}0125;\,0{,}0508]$, nên chưa đủ bằng chứng khẳng định BGE-M3 tốt hơn BM25-stemmed trên quần thể mục tiêu.
\begin{table}[H]\setstretch{1.05}
  \centering
  \caption{MRR@5 và nDCG@5 theo retriever và reranker trên tập development.}
  \label{tab:phase1-reranker-matrix}
  \small
  \setlength{\tabcolsep}{4.5pt}
  \begin{tabular}{lrrr|rrr}
    \toprule
    & \multicolumn{3}{c|}{\textbf{MRR@5}}
    & \multicolumn{3}{c}{\textbf{nDCG@5}} \\
    \textbf{Retriever} & \textbf{Không} & \textbf{MiniLM} & \textbf{BGE-large}
    & \textbf{Không} & \textbf{MiniLM} & \textbf{BGE-large} \\
    \midrule
    BGE-M3 sparse & 0,8059 & 0,8387 & \textbf{0,8797} & 0,8317 & 0,8642 & \textbf{0,8976} \\
    e5-base dense & 0,6316 & 0,7783 & 0,7997 & 0,6684 & 0,7994 & 0,8172 \\
    Hybrid RRF    & 0,7192 & 0,7938 & 0,8203 & 0,7474 & 0,8164 & 0,8405 \\
    \bottomrule
  \end{tabular}
\end{table}

\begin{table}[H]\setstretch{1.05}
  \centering
  \caption{Độ trễ P50 của ma trận truy xuất Giai đoạn 1 trên cùng môi trường chạy.}
  \label{tab:phase1-reranker-latency}
  \small\setlength{\tabcolsep}{6pt}
  \begin{tabular}{lrrr}
    \toprule
    \textbf{Retriever} & \textbf{Không rerank} & \textbf{MiniLM} &
    \textbf{BGE-large} \\
    \midrule
    BGE-M3 sparse & 66,1 ms & 163,9 ms & 510,1 ms \\
    e5-base dense & 17,7 ms & 127,3 ms & 540,5 ms \\
    Hybrid RRF    & 90,4 ms & 193,3 ms & 609,8 ms \\
    \bottomrule
  \end{tabular}
\end{table}

\begin{itemize}
\item \textbf{Reranker:} với BGE-M3, MiniLM nâng MRR@5 từ 0,8059 lên 0,8387; BGE-reranker-large đạt 0,8797. Đổi lại, P50 tăng từ 163,9 ms với MiniLM lên 510,1 ms với BGE-large.
\item \textbf{Hybrid RRF:} không vượt BGE-M3 ở cả ba chế độ xếp hạng lại; phép hợp nhất chưa đem lại lợi ích trong thiết lập khảo sát.
\end{itemize}

\noindent\textbf{Kết luận.} BGE-M3 top 20 kết hợp BGE-reranker-large top 5 được chọn cho tầng sinh, với Hit@5/Recall@5 trên development là 0,9573/0,9555.
